\subsection{Inclusion-Exclusion, Pigeonhole, and Combinatorial Proof}

\begin{definition} \label{def:multisets-and-multiplicity} \hlt{Multisets}\\
Let $S$ be a set. A \hlt{multiset on $S$} is a pair $(S_0,m)$ such that
\begin{enumerate}[label=(\roman*)]
\setlength{\itemsep}{0pt}
\item $S_0\subseteq S$ is the \hlt{underlying set};
\item $m:S_0\to\Z_{>0}$ is the \hlt{multiplicity function};
\item $m(s)$ is the \hlt{multiplicity} of $s\in S_0$;
\item if $S_0$ is finite, the \hlt{size} is $\lvert m\rvert\defeq\sum_{s\in S_0}m(s)$.
\end{enumerate}
Equivalently, record $(S_0,m)$ as $\{(s,m(s)):s\in S_0\}$.
\end{definition}

\begin{example} \label{ex:multiset-eigenvalue-bookkeeping} \hlt{Multiset Bookkeeping}\\
The list $\{a,a,b,b,b,c\}$ is the multiset on $\{a,b,c\}$ with multiplicities $2,3,1$.
This language is useful whenever order is irrelevant but repetition matters, such as eigenvalues counted with algebraic multiplicity.
\end{example}

\begin{proposition} \label{prop:finite-inclusion-exclusion-two-three-sets} \hlt{Finite Inclusion-Exclusion for Two and Three Sets}\\
Let $A,B,C$ be finite sets, and write $\mathbf{1}_E(x)$ for membership in $E$.
\begin{enumerate}[label=(\roman*)]
\setlength{\itemsep}{0pt}
\item $\abs{A\cup B}=\abs{A}+\abs{B}-\abs{A\cap B}$.
\item $\abs{A\cup B\cup C}=\abs{A}+\abs{B}+\abs{C}-\abs{A\cap B}-\abs{A\cap C}-\abs{B\cap C}+\abs{A\cap B\cap C}$.
\end{enumerate}
\begin{proof}
\begin{enumerate}[label=(\roman*)]
\setlength{\itemsep}{0pt}
\item Let $U=A\cup B$. For $x\in U$, $\mathbf{1}_U(x)=\mathbf{1}_A(x)+\mathbf{1}_B(x)-\mathbf{1}_{A\cap B}(x)$. Sum over $U$; $\abs{U}=\abs{A}+\abs{B}-\abs{A\cap B}$.
\item Let $U=A\cup B\cup C$ and $m(x)\defeq\mathbf{1}_A(x)+\mathbf{1}_B(x)+\mathbf{1}_C(x)$.\\
For $x\in U$, $m(x)\in\{1,2,3\}$ and $m(x)-\binom{m(x)}{2}+\binom{m(x)}{3}=1$.\\
Thus each $x\in U$ contributes once to $\abs{A}+\abs{B}+\abs{C}-\abs{A\cap B}-\abs{A\cap C}-\abs{B\cap C}+\abs{A\cap B\cap C}$, QED.
\end{enumerate}
\hlt{Structural Commentary}: Inclusion-exclusion weights each membership pattern to total weight $1$.\\
\hlt{Key Insight}: Add single memberships, subtract overcounts, then restore triple overlaps.
\end{proof}
\end{proposition}

\begin{definition} \label{def:symmetric-difference-of-events-and-sets} \hlt{Symmetric Difference}\\
For sets $A$ and $B$, the \hlt{symmetric difference} is $A\triangle B\defeq (A\setminus B)\cup(B\setminus A)=(A\cup B)\setminus(A\cap B)$.\\
It records the elements belonging to exactly one of $A$ and $B$.
\end{definition}

\begin{method} \label{meth:combinatorial-proof-by-classification} \hlt{Combinatorial Proof by Classification}\\
To prove $N_L=N_R$ for finite counts:
\begin{enumerate}[label=(\roman*)]
\setlength{\itemsep}{0pt}
\item choose a finite object set $\Omega$ and a pattern map $p:\Omega\to P$;
\item write both sides as sums over fibers, $N_L=\sum_{\alpha\in P}w_L(\alpha)\abs{p^{-1}(\{\alpha\})}$ and $N_R=\sum_{\alpha\in P}w_R(\alpha)\abs{p^{-1}(\{\alpha\})}$;
\item verify $w_L(\alpha)=w_R(\alpha)$ for every $\alpha\in P$.
\end{enumerate}
\end{method}
